\section{The ordered term structure, and the contract}
\label{sec:vsterm}

Everything so far concerned one expiry.  A surface carries a fair
strike per expiry, and the last question is how the family of numbers
$w_{\rm vs}(\tau)$ behaves across maturities.  The answer is that the
calendar condition of Chapter~2 orders it --- and the proof is a
three-line reuse of the spanning identity.

Recall the condition.  A surface is calendar-ordered when, at every
log-moneyness, the later expiry's normalized call is worth at least
the earlier one's: $c_1(k)\le c_2(k)$ for all $k$
(\cref{eq:calreq}).  Chapter~2 showed this is the same as
\emph{convex order} of the two laws --- every convex payoff is priced
at least as high under the later law --- and, by Kellerer's theorem,
exactly the condition for some martingale to have these two laws as
its distributions at the two expiries
(\cref{thm:ledgerorder,sec:calendar}).

\begin{proposition}[Calendar order integrates]
\label{prop:vsterm}
If two admissible mean-one slices satisfy $c_1(k)\le c_2(k)$ for all
$k$, then $w_{\rm vs,1}\le w_{\rm vs,2}$.
\end{proposition}

\begin{proof}
Puts inherit the order from calls: at a common strike $K$, parity
gives $P=c+K-1$ for both slices, so $c_1\le c_2$ implies
$P_1\le P_2$.  Now write each fair strike through
\cref{prop:vsspanning} with $\varphi(y)=-2\log y$, whose second
derivative $\varphi''(K)=2/K^{2}$ is positive:
\[
w_{\rm vs}
=2\int_0^{1}\frac{P(K)}{K^{2}}\,\dd K
+2\int_1^{\infty}\frac{c(K)}{K^{2}}\,\dd K .
\]
Every integrand of slice~2 dominates the corresponding integrand of
slice~1, so the integrals are ordered.
\end{proof}

The same argument proves more: \emph{any} payoff with
$\varphi''\ge0$ --- any convex payoff --- has its price ordered by
the calendar, which is just convex order saying its own name.  The
variance swap is the special case a desk trades directly, and for it
the order has a market reading.  The difference
$w_{\rm vs}(\tau_2)-w_{\rm vs}(\tau_1)$ is the \emph{forward
variance} between the two dates --- the variance a client buys by
holding the long swap against the short one.  \Cref{prop:vsterm}
says a calendar-ordered surface never prices forward variance below
zero; a violation would let a desk lock in variance at a negative
price, which is the calendar arbitrage of Chapter~2 restated for the
integral.

\begin{figure}[htbp]
\centering
\paperfig{\textwidth}{fig_vs_term.pdf}
\caption{The integral's term structure on the frozen gallery (stored
surface fits).  (a)~SPY in total variance: the fair strike
$w_{\rm vs}$ and the ATM total variance $w(0)$ both rise with
maturity, and the gap between them widens.  (b)~Both names in
volatility units: the fair-strike vol rides above the ATM vol
everywhere, by up to \MacVsTermSpySpreadBp\ vol bp for SPY and
\MacVsTermNvdaSpreadBp\ for NVDA.  NVDA's vol \emph{levels} are not
monotone in maturity --- the order lives in total variance, panel
(a)'s units, not in vol.}
\label{fig:vsterm}
\end{figure}

\Cref{fig:vsterm} reads the term structure off the frozen gallery ---
the snapshot's stored surface fits, the same objects as
\cref{fig:vsshare}(b).  In panel~(a), SPY in total variance: the
orange curve is $w_{\rm vs}(\tau)$, the blue is the ATM total
variance $w(0)(\tau)$, both increasing.  Across every adjacent pair
of expiries, in both names, the fair strike's increment is positive
--- the smallest measured increment is $+\MacVsTermIncMinVarBp$
variance bp (a variance bp is $10^{-4}$ of total variance, the
sibling of the vol bp of \cref{sec:vsnumber}), between the two front
expiries a day apart, and the smallest forward variance, expressed as
a forward volatility, is
\MacVsTermFwdVolMinPct\% --- so the published surface keeps the order
\cref{prop:vsterm} demands, with its thinnest margin exactly where
maturities crowd together.  Panel~(a) shows something else: the
\emph{gap} between the curves widens with maturity.  By the
flat-smile computation \eqref{eq:vsflat}, that gap is precisely the
part of the fair strike owed to the smile's departure from flatness,
and it compounds with maturity just as the unquoted share of
\cref{fig:vsshare}(b) did --- the two figures are the same fact in
two currencies.

Panel~(b) restates the term structure in the units desks quote,
volatility, for both names.  The fair-strike vol $\sigma_{\rm vs}$
rides above the ATM vol at every expiry --- by up to
\MacVsTermSpySpreadBp\ vol bp at SPY's long end and
\MacVsTermNvdaSpreadBp\ for NVDA.  On a flat smile the two curves
would coincide; the spread \emph{is} the priced skew and wing mass,
and a desk can read it directly as what convexity in the smile costs.
One trap is visible and worth flagging: NVDA's vol curves are not
monotone in maturity, and nothing is wrong --- the calendar orders
\emph{total variance}, panel~(a)'s units, and a falling vol term
structure is perfectly consistent with rising total variance.
Ordering vols is not a no-arbitrage condition; ordering their
$\sigma^{2}\tau$ is.

\subsection{The contract}
\label{sec:vscontract}

The chapter closes the way the model chapters do, with its contract
--- three columns of decreasing strength.  The first holds by
theorem.  The second is measured on the book's frozen data, and the
chapter quoted every number.  The third is convention: chosen, stated,
and owed to the user, because no data selects it.

\begin{table}[htbp]
\centering
\small
\caption{The chapter's contract.}
\label{tab:vscontract}
\begin{tabular}{p{0.31\textwidth}p{0.31\textwidth}p{0.30\textwidth}}
\toprule
\textbf{Proved} & \textbf{Checked (frozen data)} & \textbf{Convention}\\
\midrule
The identity chain \eqref{eq:vsmaster}: one fair strike, three exact
integrals (smile, law, field). &
Law vs strike route on one slice: \MacVsAgreeRankBp\ vol bp; field vs
strike route on one sheet: \MacVsAgreeFieldCoarseBp\ vol bp at the
march's step, \MacVsAgreeFieldBp\ at an eightfold refinement. &
Which route a model evaluates (its native chart), and the reading of
the fair strike: diffusion paths, continuous monitoring, variance
time.  Jumps sit outside the identity
\eqref{eq:vslogcontract}.\\
\addlinespace
The ceiling (\cref{prop:vsceiling}): no admissible wing steeper than
$2$; the boundary ray itself inadmissible. &
Fitted slopes (\cref{tab:vswings}): all wings far inside the
ceiling. &
The wing drawn \emph{within} the cone: each family's stated
mechanism --- structural, capped, or inherited-and-diagnosed.\\
\addlinespace
The envelope (\cref{prop:vscone}): every completion between the
convex floor and the monotone cap; top edge climbs at the ceiling
\eqref{eq:vstopslope}. &
All three families' wings clear their cones (clearances of
\cref{sec:vsenvelope}); the hatted counterexample caught by the
extended $g_{\rm D}$ check, slopes unmoved
(\cref{fig:winglimit}). &
Where conditions are enforced: one-curve checks extended into the
wings; two-curve (calendar) checks extended under the stack's stated
tail policy (\cref{eq:globalledgerconstraint}).\\
\addlinespace
Calendar order integrates (\cref{prop:vsterm}): the fair strike's
term structure is nondecreasing; forward variance is nonnegative. &
Smallest measured increment across both names:
$+\MacVsTermIncMinVarBp$ variance bp; smallest forward vol
\MacVsTermFwdVolMinPct\%. &
The weight given a traded var-swap quote when one exists: one more
residual in the objective of \cref{sec:calibration}, at a declared
strength.\\
\bottomrule
\end{tabular}
\end{table}

Step back, once, to the book's thesis.  The fair strike is a single
number in which the three ingredients the thesis separates are
unusually easy to see.  Validity supplied the theorems: the identity
chain, the ceiling, the envelope, the order.  Information is the
quoted belly's \MacVsShareQuotedPct\%.  Convention is the rest ---
the wing contract that completes the integral, stated in a table
rather than hidden in an extrapolation routine.  Part~I is now
complete: laws, smiles, fields, and the integrals that test them.
Everything so far took the forward and the discount for granted at
every expiry.  Part~II starts by inferring them from the quotes.
